% Part: first-order-logic
% Chapter: model-theory
% Section: theory-of-m

\documentclass[../../../include/open-logic-section]{subfiles}

\begin{document}

\olfileid{mod}{bas}{thm}

\section{Alles wat in een \printtoken{S}{structure} waar is: haar theorie}

In elke !!{structure}~$\Struct{M}$ zijn sommige !!{sentence}s waar en
andere niet. Neem alle !!{sentence}s die erin waar zijn bij elkaar. Die
verzameling heet de \emph{theorie} van de structuur. Het is ook echt een
theorie. Alles wat uit die verzameling volgt, is namelijk waar in elk
model ervan, en dus ook in~$\Struct{M}$ zelf.

\begin{defn}
  Neem !!a{structure}~$\Struct M$. De \emph{theorie} van
  $\Struct{M}$ is de verzameling $\Theory{M}$ van alle !!{sentence}s
  die in $\Struct{M}$ waar zijn. Met symbolen: $\Theory{M} =
  \Setabs{!A}{\Sat{M}{!A}}$.
\end{defn}

We gebruiken het woord ``theorie'' ook losser. Dan bedoelen we een
verzameling !!{sentence}s met een bepaalde bedoelde uitleg. Het maakt
voor dat losse gebruik niet uit of alles wat eruit volgt al in de
verzameling staat, dus of zij deductief gesloten is.

\begin{prop}
Voor elke $\Struct{M}$ is $\Theory{M}$ compleet.
\end{prop}

\begin{proof}
Neem een willekeurige !!{sentence}~$!A$. Dan geldt ofwel
$\Sat{M}{!A}$, ofwel $\Sat{M}{\lnot !A}$. Daarom staat ofwel
$!A \in \Theory{M}$, ofwel $\lnot !A \in \Theory{M}$. De theorie kiest
dus voor elke zin de zin zelf of haar ontkenning.
\end{proof}

\begin{prop}\ollabel{prop:equiv}
  Als $\Struct{N} \models !A$ voor elke $!A \in \Theory{M}$, dan
  $\Struct{M} \elemequiv \Struct{N}$.
\end{prop}

\begin{proof}
Uit $\Sat{N}{!A}$ voor alle $!A \in \Theory{M}$ volgt
$\Theory{M} \subseteq \Theory{N}$. Stel nu dat $\Sat{N}{!A}$. Dan
$\Sat/{N}{\lnot !A}$, zodat $\lnot !A \notin \Theory{M}$. De verzameling
$\Theory{M}$ is compleet. Daarom moet $!A \in \Theory{M}$. We hebben dus
$\Theory{N} \subseteq \Theory{M}$, en daarmee
$\Struct{M} \elemequiv \Struct{N}$.
\end{proof}

\begin{rem}\ollabel{remark:R}
  Bekijk $\Struct{R} = \langle\Real, <\rangle$. Dit is de !!{structure}
  waarvan het domein de verzameling $\Real$ van alle reële getallen is.
  De !!{language} heeft maar één 2-plaatsig !!{predicate}; dat betekent
  in deze structuur de relatie $<$ op de reële getallen. De structuur
  $\Struct{R}$ is duidelijk !!{nonenumerable}. Toch is $\Theory{R}$
  consistent. Volgens de stelling van L\"owenheim--Skolem heeft die
  theorie daarom !!a{enumerable} model; noem dat model $\Struct{S}$.
  Uit \olref{prop:equiv} volgt dan $\Struct{R} \equiv \Struct{S}$.
  Maar $\Struct{R}$ en $\Struct{S}$ zijn niet isomorf. De omgekeerde
  richting van \olref[iso]{thm:isom} geldt dus niet altijd.
\end{rem}

\end{document}
